\scanpage{68}
% upper table, counts in parentheses are placed beneath the triples
\uline{Cas tétraédraux}
\diagram{$\begin{array}{c@{\quad}lll}
&\underset{(4)}{(3,3,2)}&\underset{(2)}{\underline{(3,2,3)}}&\underset{(2)}{\underline{(2,3,3)}}\\
6&\underset{(4)}{(3,3,3)}
\end{array}$}

\uline{Cas octaédraux}
\diagram{\scriptsize\setlength{\arraycolsep}{2pt}$\begin{array}{c@{\quad}llllll}
&\underset{(4)}{(4,3,2)}&\underset{(2)}{\underline{(3,2,4)}}&\underset{(2)}{\underline{(2,4,3)}}&
\underset{(2)}{\underline{(2,3,4)}}&\underset{(4)}{(3,4,2)}&\underset{(2)}{\underline{(4,2,3)}}\\
13&\underset{(4)}{(4,4,3)}&\underset{(4)}{(4,\deleted{\unsure{4}}3,4)}&\underset{(4)}{(3,4,4)}
\end{array}$}

\uline{Cas icosaédraux}
\diagram{\scriptsize\setlength{\arraycolsep}{2pt}$\begin{array}{r@{\quad}l}
\begin{array}{r}10\\+4\end{array}&
\left\{\begin{array}{llllll}
\underset{(8)}{(5,\deleted{\unsure{5}}3,2)}&
\underset{(2)}{\underline{(3,2,5)}}&\underset{(4)}{\underline{(2,5,3)}}&
\underset{(2)}{\underline{(2,3,5)}}&\underset{(8)}{(3,5,2)}&
\underset{(4)}{\underline{(5,2,3)}}\\
&\underset{(2)}{\underline{(3,2,5')}}&\deleted{\unsure{(}}&
\underset{(2)}{\underline{(2,3,5')}}
\end{array}\right.\\[5pt]
\begin{array}{r}+4\\2\end{array}&
\left\{\begin{array}{lll}
\underset{(16)}{(5,5,3)}&\underset{(8)}{(5,3,5)}&\underset{(8)}{(3,5,5)}\\
\underset{(8)}{(\deleted{\unsure{3}}5,\deleted{\unsure{5}}3,5')}&
\underset{(8)}{(\deleted{\unsure{5}}3,\deleted{\unsure{3}}5,5')}
\end{array}\right.\\[5pt]
\begin{array}{r}4\\1\\\hline25\end{array}&
\left\{\begin{array}{ll}
\underset{(4)}{\underline{(2,5,5')}}&\underset{(4)}{\underline{(5,2,5')}}\\
\underset{(16)}{(5,5,5)}
\end{array}\right.
\end{array}$}

% right-hand note begins beside the second icosahedral group
\margin{\uline{NB} On a noté en dessous de chaque symbole le nombre d'hyperpolygônes non normalisés qu'on peut \deleted{en} déduire de chaque repr. normalisé, par le procédé $\tau_i\mapsto\tau_i^{k_i}$ [égal à $\varphi(\nu_0)\varphi(\nu_1)$\unsure{,|;} prend les valeurs $2,4,8,16$]}

% summary in three columns
\diagram{$\begin{array}{rl}
6&\text{Cas tétraédraux}\\13&\text{Cas octaédraux}\\
\text{\uuline{25}}&\text{Cas icosaédraux}\\44&\text{Cas au total}
\end{array}$}
\margin{[Il y a 14 cas sur les 44 où le type numérique ne suffit pas à déterminer le type d'iso\unsure{morphie|morphi.}]}
\margin{On vérifie bien que le total (\unsure{des|obtenu} \deleted{\unsure{xx}} \unsure{des} pour \unsure{l'ens|l'un|l'ens.} de tous les cas) redonne $4\cdot45=180$. Dans le calcul, il faut bien sûr tenir compte aussi des soulignés, pour compter double}

% final paragraph, continues on the following scan
[ \uline{NB} Les syst. pythagoriciens de 2 ps. réflexions $\tau_0,\tau_{\unsure{1|4}}$ (sans normalisation sur $\lambda_0,\lambda_1$) sont plus nombreux, chacun des 45 systèmes de \unsure{Schwartz|Schwarz} $(K_0,K_1,K)$ possibles donnant naissance à 4 \inserted{$=2\times2$} \uline{possibilités}, d'où 180 cas possibilités. Pour calculer les $\zeta$ dans chacun de ces 180 cas, le travail est simplifié par la remarque que $\zeta$ ne
% archive sheet number: 67
